\documentclass[../../../include/open-logic-section]{subfiles}

\begin{document}

\olfileid[hi]{sth}{ordinals}{replacement}
\olsection{प्रतिस्थापन}

\olref[sth][ordinals][vn]{sec} में हमने यह सुझाव देकर क्रमसूचकों को प्रस्तुत
करने की प्रेरणा दी थी कि उन्हें क्रम-प्रकार, अर्थात सुव्यवस्थाओं के विहित
प्रतिनिधि, माना जा सकता है। यह सफल होने के लिए हमें सिद्ध करना होगा कि
\emph{प्रत्येक सुव्यवस्था किसी क्रमसूचक के समाकृतिक है}। इससे हम
$\ordtype{A, <}$ को ऐसा क्रमसूचक $\alpha$ परिभाषित कर सकेंगे कि
$\tuple{A, <} \isomorphic \alpha$।

दुर्भाग्य से अब तक प्रस्तुत अभिगृहीतों से ही हम वांछित परिणाम सिद्ध
\emph{नहीं कर सकते}। (क्यों, यह \olref[replacement][strength]{sec} में
देखेंगे; अभी बात केवल इतनी है कि हम नहीं कर सकते।) हमें एक नया विचार चाहिए,
और वह यह है:

\begin{axiom}[प्रतिस्थापन की योजना]
किसी भी सूत्र $\phi(x, y)$ के लिए निम्न एक अभिगृहीत है:
\begin{quote}
	किसी भी $A$ के लिए यदि $(\forall x \in A)\lexists![y][\phi(x,y)]$,
	तो $\Setabs{y}{(\exists x \in A)\phi(x,y)}$ का अस्तित्व है।
\end{quote}
\end{axiom}
\noindent
पृथक्करण की तरह यह एक योजना है: यह अनंततः अनेक भिन्न $\phi$ में प्रत्येक
के लिए एक, इस प्रकार अनंततः अनेक अभिगृहीत देती है। इसे समान रूप से (और
सामान्यतः) इस तरह लिखा जाता है:

\begin{defish}
ऐसे किसी भी सूत्र $\phi(x,y)$ के लिए जिसमें ``$B$'' न हो, निम्न एक
अभिगृहीत है:
\[
\forall A[(\forall x \in A)\lexists![y][\phi(x,y)] \lif \exists B\forall y (y \in B \liff (\exists x \in A)\phi(x,y))]
\]
\end{defish}

पहली बार देखने पर यह काफ़ी उलझा हुआ सूत्र है। प्रतिस्थापन का निम्न त्वरित
परिणाम संभवतः हमारे सहज विचार को \emph{अधिक स्पष्ट} रूप में व्यक्त करता
है:\footnote{पद ऐसा व्यंजक है जो (अपने मुक्त चरों के किसी भी पूरण पर) ठीक
एक वस्तु चुनता है। उदाहरणतः ``$\emptyset$'' रिक्त समुच्चय चुनने वाला पद है;
``$\{x\}$'' $x$ का एकल समुच्चय चुनने वाला पद है ($x$ जो भी हो); और
``$x \cup y$'' $x$ और $y$ का सम्मिलन चुनने वाला पद है (वे जो भी हों)।}

\begin{cor}
किसी भी पद $\tau(x)$ और किसी भी समुच्चय $A$ के लिए इस समुच्चय का अस्तित्व है:
\[
	\Setabs{\tau(x)}{x \in A} = \Setabs{y}{(\exists x \in A)y = \tau(x)}.
\]
\end{cor}

\begin{proof}
चूँकि $\tau$ एक \emph{पद} है,
$\forall x \lexists![y][\tau(x) = y]$। अतः और भी अधिक स्पष्टतः
$(\forall x \in A)\lexists![y][\tau(x) = y]$। इसलिए प्रतिस्थापन से
$\Setabs{y}{(\exists x \in A)\tau(x) = y}$ का अस्तित्व है।
\end{proof}
\noindent
यह सुझाता है कि अभिगृहीत के लिए ``प्रतिस्थापन'' अच्छा नाम है: कोई समुच्चय
$A$ दिए जाने पर आप $A$ के प्रत्येक सदस्य को~$\tau$ के अधीन उसके प्रतिबिंब
से बदलकर नया समुच्चय $\Setabs{\tau(x)}{x \in A}$ बना सकते हैं। वास्तव में
किसी फलन के अधीन समुच्चय के प्रतिबिंब के संकेतन का अनुसरण करते हुए हम
$\Setabs{\tau(x)}{x \in A}$ के लिए $\funimage{\tau}{A}$ लिख सकते हैं।

किंतु निर्णायक बात यह है कि $\tau$ एक \emph{पद} है।
\olref[sfr][fun][rel]{sec} के अर्थ में उसका कोई \emph{फलन} (का नाम), अर्थात
क्रमित युग्मों का कोई विशिष्ट समुच्चय, होना आवश्यक नहीं। आखिर यदि $f$ (उस
अर्थ में) फलन है, तो समुच्चय
$\funimage{f}{A} = \Setabs{f(x)}{x \in A}$ केवल $\ran{f}$ का एक विशिष्ट
उपसमुच्चय है, और केवल~$\Zminus$ के अभिगृहीतों से उसके अस्तित्व की गारंटी
पहले ही है।\footnote{बस
$\Setabs{y \in \bigcup \bigcup f}{(\exists x \in A)y = f(x)}$ पर विचार
करें।} इसके विपरीत प्रतिस्थापन हमारे अभिगृहीतों में एक \emph{शक्तिशाली}
योग है, जैसा हम \olref[sth][replacement][]{chap} में देखेंगे।


\end{document}
